% !Mode:: "TeX:UTF-8"
% !TEX root = ../../main.tex
% ============================================================
% TOPIC 9
% ============================================================
\mashgulot{9}{The Gauss and Gauss–Jordan methods for systems of linear equations}

\begin{tarif}
\textbf{Gaussian elimination} is a general method for solving systems of
linear equations. Elementary row operations are applied to the augmented
matrix of the system in order to bring it to \textbf{row echelon form}; the
unknowns are then found by \textbf{back substitution}.

Accordingly, the method consists of two stages. In the \textbf{forward
elimination} the augmented matrix is reduced to row echelon form; in the
\textbf{back substitution} the unknowns are determined one after another,
starting from the last equation.
\end{tarif}

\begin{tarif}
The \textbf{Gauss–Jordan method} is a modification of Gaussian elimination
in which the reduction is continued beyond row echelon form until the matrix
reaches \textbf{reduced row echelon form} (RREF). No back substitution is
then needed: the values of the unknowns are read off directly.
\end{tarif}

For a system of $m$ equations in $n$ unknowns the \textbf{augmented matrix} is
\[\kf{(A\,|\,b) = \left(\begin{array}{cccc|c}
a_{11}&a_{12}&\cdots&a_{1n}&b_1\\
a_{21}&a_{22}&\cdots&a_{2n}&b_2\\
\vdots&\vdots&\ddots&\vdots&\vdots\\
a_{m1}&a_{m2}&\cdots&a_{mn}&b_m
\end{array}\right).}\]

Unlike Cramer’s rule and the matrix method, Gaussian elimination works also
when the number of equations differs from the number of unknowns
($m\neq n$) and when $\det A = 0$. In the row operations below, $S_i$ denotes
the $i$th row of the matrix.

\subsection{Gaussian elimination: a system with a unique solution}

\begin{misol}
$\begin{cases}
2x_1 + x_2 - x_3 = 1,\\
x_1 - x_2 + 2x_3 = 5,\\
3x_1 + 2x_2 + x_3 = 8.
\end{cases}$

\textbf{Solution.}
We begin the forward elimination with the augmented matrix. To have a
leading coefficient equal to 1, we first interchange the first two rows:
\[\kf{\left(\begin{array}{ccc|c}
2&1&-1&1\\1&-1&2&5\\3&2&1&8
\end{array}\right)
\xrightarrow{S_1\leftrightarrow S_2}
\left(\begin{array}{ccc|c}
1&-1&2&5\\2&1&-1&1\\3&2&1&8
\end{array}\right).}\]
Subtracting suitable multiples of the new first row from the other two rows
and then eliminating the second entry of the third row, we obtain
\[\kf{\xrightarrow[\;S_3-3S_1\;]{S_2-2S_1}
\left(\begin{array}{ccc|c}
1&-1&2&5\\0&3&-5&-9\\0&5&-5&-7
\end{array}\right)
\xrightarrow{S_3-\frac{5}{3}S_2}
\left(\begin{array}{ccc|c}
1&-1&2&5\\0&3&-5&-9\\0&0&\frac{10}{3}&8
\end{array}\right).}\]
The back substitution starts with the last equation, which gives
\[\kf{\frac{10}{3}x_3 = 8 \;\Rightarrow\; x_3 = \frac{24}{10} = \frac{12}{5}.}\]
Substituting this value into the second equation, we find
\[\kf{3x_2 - 5\cdot\frac{12}{5} = -9 \;\Rightarrow\;
3x_2 = -9+12 = 3 \;\Rightarrow\; x_2 = 1,}\]
and the first equation then yields
\[\kf{x_1 - 1 + 2\cdot\frac{12}{5} = 5 \;\Rightarrow\;
x_1 = 5+1-\frac{24}{5} = \frac{30-24}{5} = \frac{6}{5}.}\]

The fractions here are a little inconvenient. In hand computations one
often rescales a row, for instance $S_3\to\frac{3}{10}S_3$, to keep the
numbers simple; the result, of course, does not depend on this.

\textbf{Answer:} $x_1=\frac{6}{5}$, $x_2=1$, $x_3=\frac{12}{5}$.
\end{misol}

\begin{misol}
$\begin{cases}
x_1 + x_2 + x_3 = 6,\\
x_1 - x_2 + 2x_3 = 5,\\
2x_1 + x_2 - x_3 = 1.
\end{cases}$

\textbf{Solution.}
Here the first coefficient is already 1, so the forward elimination
proceeds at once:
\[\kf{\left(\begin{array}{ccc|c}
1&1&1&6\\1&-1&2&5\\2&1&-1&1
\end{array}\right)
\xrightarrow[\;S_3-2S_1\;]{S_2-S_1}
\left(\begin{array}{ccc|c}
1&1&1&6\\0&-2&1&-1\\0&-1&-3&-11
\end{array}\right)}\]
\[\kf{\xrightarrow{S_3-\frac{1}{2}S_2}
\left(\begin{array}{ccc|c}
1&1&1&6\\0&-2&1&-1\\[3pt]
0&0&-\frac{7}{2}&-\frac{21}{2}
\end{array}\right).}\]
Back substitution gives successively
\[\kf{-\frac{7}{2}x_3 = -\frac{21}{2} \;\Rightarrow\; x_3 = 3,\qquad
-2x_2 + 3 = -1 \;\Rightarrow\; x_2 = 2,\qquad
x_1 + 2 + 3 = 6 \;\Rightarrow\; x_1 = 1.}\]

\textbf{Answer:} $x_1=1$, $x_2=2$, $x_3=3$.

Check: $1+2+3=6$, \ $1-2+6=5$, \ $2+2-3=1$.
\end{misol}

\subsection{The Gauss–Jordan method}

In the Gauss–Jordan method the matrix is brought to reduced row echelon
form: every leading entry equals 1, and all other entries in the column of a
leading entry, above and below it, are zero.

\begin{misol}
Let us solve the system of the previous example by the Gauss–Jordan method,
starting from the row echelon form obtained there. We first scale the third
row so that its leading entry becomes 1:
\[\kf{\left(\begin{array}{ccc|c}
1&1&1&6\\0&-2&1&-1\\[3pt]
0&0&-\frac{7}{2}&-\frac{21}{2}
\end{array}\right)
\xrightarrow{S_3\to -\frac{2}{7}S_3}
\left(\begin{array}{ccc|c}
1&1&1&6\\0&-2&1&-1\\
0&0&1&3
\end{array}\right).}\]
Next we clear the entry above it in the second row and normalize the second
row:
\[\kf{\xrightarrow{S_2-S_3}
\left(\begin{array}{ccc|c}
1&1&1&6\\0&-2&0&-4\\
0&0&1&3
\end{array}\right)
\xrightarrow{S_2\to -\frac{1}{2}S_2}
\left(\begin{array}{ccc|c}
1&1&1&6\\0&1&0&2\\
0&0&1&3
\end{array}\right).}\]
Finally, subtracting the second and third rows from the first, we arrive at
\[\kf{\xrightarrow[\;S_1-S_3\;]{S_1-S_2}
\left(\begin{array}{ccc|c}
1&0&0&1\\0&1&0&2\\
0&0&1&3
\end{array}\right).}\]

The solution can now be read off directly: $x_1=1$, $x_2=2$, $x_3=3$.
\end{misol}

\subsection{Inconsistent systems and systems with infinitely many solutions}

\textbf{An inconsistent system.}

\begin{misol}
$\begin{cases}
x_1 + 2x_2 = 3,\\
2x_1 + 4x_2 = 5.
\end{cases}$

\textbf{Solution.}
\[\kf{\left(\begin{array}{cc|c}
1&2&3\\2&4&5
\end{array}\right)
\xrightarrow{S_2-2S_1}
\left(\begin{array}{cc|c}
1&2&3\\0&0&-1
\end{array}\right).}\]
The last row corresponds to the equation $0\cdot x_1+0\cdot x_2=-1$, which
is a contradiction. Hence the system is \textbf{inconsistent}: it has no
solution.
\end{misol}

\textbf{A system with infinitely many solutions.}

\begin{misol}
$\begin{cases}
x_1 - x_2 + 2x_3 = 3,\\
2x_1 - 2x_2 + 4x_3 = 6,\\
x_1 + x_2 - x_3 = 1.
\end{cases}$

\textbf{Solution.}
\[\kf{\left(\begin{array}{ccc|c}
1&-1&2&3\\2&-2&4&6\\1&1&-1&1
\end{array}\right)
\xrightarrow[\;S_3-S_1\;]{S_2-2S_1}
\left(\begin{array}{ccc|c}
1&-1&2&3\\0&0&0&0\\0&2&-3&-2
\end{array}\right)
\xrightarrow{S_2\leftrightarrow S_3}
\left(\begin{array}{ccc|c}
1&-1&2&3\\0&2&-3&-2\\0&0&0&0
\end{array}\right).}\]
The ranks of the coefficient matrix and of the augmented matrix coincide,
$\operatorname{rank}A = \operatorname{rank}(A\,|\,b) = 2$, but they are less
than the number of unknowns, which is 3. So the system is consistent and has
one free unknown; we take $x_3 = t$ as a parameter. The second row gives
$2x_2-3t=-2$, that is, $x_2=\frac{3t-2}{2}$, and from the first row
\[\kf{x_1 = 3+x_2-2t = 3+\frac{3t-2}{2}-2t = 2-\frac{t}{2}.}\]

\textbf{Answer:} $x_1=2-\frac{t}{2}$, $x_2=\frac{3t-2}{2}$, $x_3=t$, where
$t\in\mathbb{R}$ is arbitrary.
\end{misol}

\begin{misol}
$\begin{cases}
x_1 + 2x_2 - x_3 + x_4 = 2,\\
2x_1 + 4x_2 - 2x_3 + 3x_4 = 5,\\
x_1 + 2x_2 - x_3 + 2x_4 = 3.
\end{cases}$

\textbf{Solution.}
\[\kf{\left(\begin{array}{cccc|c}
1&2&-1&1&2\\2&4&-2&3&5\\1&2&-1&2&3
\end{array}\right)
\xrightarrow[\;S_3-S_1\;]{S_2-2S_1}
\left(\begin{array}{cccc|c}
1&2&-1&1&2\\0&0&0&1&1\\0&0&0&1&1
\end{array}\right)
\xrightarrow{S_3-S_2}
\left(\begin{array}{cccc|c}
1&2&-1&1&2\\0&0&0&1&1\\0&0&0&0&0
\end{array}\right).}\]
Here $\operatorname{rank}A = \operatorname{rank}(A\,|\,b) = 2$ and there are
four unknowns, so $4-2=2$ of them are free. We choose $x_2=s$ and $x_3=t$.
The second row gives $x_4 = 1$, and then the first row gives
$x_1 = 2-2s+t-1 = 1-2s+t$.

\textbf{Answer:} $x_1=1-2s+t$, $x_2=s$, $x_3=t$, $x_4=1$, where
$s,t\in\mathbb{R}$.
\end{misol}

\textbf{Comparison of the Gauss and Gauss–Jordan methods.}

\begin{center}
\begin{tabular}{|p{5.3cm}|p{4.3cm}|}
\hline
\textbf{Gaussian elimination} & \textbf{Gauss–Jordan method} \\
\hline
Reduction to row echelon form followed by back substitution &
Reduction to RREF; the answer is read off directly \\
\hline
Needs fewer operations &
Needs more operations, but the procedure is uniform \\
\hline
Back substitution requires solving the equations one by one &
No further solving is needed \\
\hline
Also yields the determinant and the rank &
Also yields the inverse matrix: $(A\,|\,E)\to(E\,|\,A^{-1})$ \\
\hline
\end{tabular}
\end{center}

\subsection{Applications of Gaussian elimination in information security}

\textbf{Linear equations in cryptographic algorithms.}
The security analysis of many modern ciphers (AES, linearized versions of
GOST, and others) uses the method of \textbf{linear cryptanalysis}. One
writes down linear approximations of the encryption process, which leads to
a large system of linear equations. Gaussian elimination is the main tool
for solving such systems, because they may contain thousands of equations,
they need not have a unique solution ($\det A=0$ is possible), and they
usually have to be solved over the binary field $\mathrm{GF}(2)$.

\textbf{Gaussian elimination over $\mathrm{GF}(2)$.}
In information security one often works over the field
$\mathrm{GF}(2) = \{0,1\}$, in which addition is the XOR operation
($\oplus$) and multiplication is the AND operation. Gaussian elimination
works over this field without change. Adding one row to another is a
bitwise XOR, and no scaling of rows is ever needed, since the only nonzero
element is $1$ and $1^{-1}=1$. This version of the method is widely used in
cryptanalysis, in coding theory and in the analysis of hash functions.

\begin{misol}
Solve the following system over $\mathrm{GF}(2)$:
$\begin{cases}
x_1\oplus x_2\oplus x_3=1,\\
x_1\oplus x_3=0,\\
x_2\oplus x_3=1.
\end{cases}$

\textbf{Solution.}
\[\kf{\left(\begin{array}{ccc|c}
1&1&1&1\\1&0&1&0\\0&1&1&1
\end{array}\right)
\xrightarrow{S_2\oplus S_1}
\left(\begin{array}{ccc|c}
1&1&1&1\\0&1&0&1\\0&1&1&1
\end{array}\right)
\xrightarrow{S_3\oplus S_2}
\left(\begin{array}{ccc|c}
1&1&1&1\\0&1&0&1\\0&0&1&0
\end{array}\right).}\]
Back substitution gives $x_3=0$, $x_2=1$ and
$x_1=1\oplus x_2\oplus x_3=1\oplus 1\oplus 0=0$.

\textbf{Answer:} $x_1=0$, $x_2=1$, $x_3=0$.
\end{misol}

\textbf{Syndrome decoding in coding theory.}
To detect and correct errors in a received word one has to solve the
syndrome equation $Hc^T=s$. For large codes this is done by Gaussian
elimination.

\textbf{Lattice-based cryptography.}
In LWE and other lattice-based schemes, an attack that recovers the secret
key reduces to solving a system of the form $Ax\approx b \pmod{q}$.
Gaussian elimination and its modular version are standard tools in the
security analysis of such schemes.

\subsection{Computing with Python, Maple and Mathcad}

The reduced row echelon form of an augmented matrix is computed by a single
command in each of the three systems. We take the system of the second
example. In SymPy the method \texttt{rref()} returns the reduced matrix
together with the indices of the pivot columns.
\begin{Verbatim}[fontsize=\footnotesize,frame=single,label={\small Python}]
from sympy import Matrix
Ab = Matrix([[1, 1, 1, 6],
             [1, -1, 2, 5],
             [2, 1, -1, 1]])
R, pivots = Ab.rref()
print(R)
print(pivots)
\end{Verbatim}
\begin{Verbatim}[fontsize=\footnotesize,frame=single,label={\small Output}]
Matrix([[1, 0, 0, 1], [0, 1, 0, 2], [0, 0, 1, 3]])
(0, 1, 2)
\end{Verbatim}
The last column of $R$ contains the solution $x_1=1$, $x_2=2$, $x_3=3$.

In Maple, \texttt{GaussianElimination} produces a row echelon form (the
result of the forward elimination), while \texttt{ReducedRowEchelonForm}
carries out the full Gauss–Jordan reduction.
\begin{Verbatim}[fontsize=\footnotesize,frame=single,label={\small Maple}]
with(LinearAlgebra):
A := Matrix([[1, 1, 1], [1, -1, 2], [2, 1, -1]]):
b := Vector([6, 5, 1]):
GaussianElimination(<A | b>);
ReducedRowEchelonForm(<A | b>);
# [[1, 0, 0, 1], [0, 1, 0, 2], [0, 0, 1, 3]]
\end{Verbatim}

In Mathcad Prime the augmented matrix is formed with \texttt{augment(A, b)}
and reduced with the built-in function \texttt{rref}.
\begin{Verbatim}[fontsize=\footnotesize,frame=single,label={\small Mathcad}]
     | 1   1   1 |          | 6 |
A := | 1  -1   2 |     b := | 5 |
     | 2   1  -1 |          | 1 |

                       | 1  0  0  1 |
rref(augment(A, b)) =  | 0  1  0  2 |
                       | 0  0  1  3 |
\end{Verbatim}

\subsection{Exercises on the Gauss and Gauss–Jordan methods}
\bigskip

\emph{1. Solve the following $2\times 2$ systems by Gaussian elimination.}
\bigskip
\begin{multicols}{2}\raggedright
\begin{enumerate}[label=\arabic*., leftmargin=*]

\item $\begin{cases}x_1+3x_2=7\\2x_1-x_2=0\end{cases}$

\item $\begin{cases}3x_1-2x_2=4\\x_1+x_2=3\end{cases}$

\item $\begin{cases}-x_1+4x_2=5\\2x_1+x_2=7\end{cases}$

\item $\begin{cases}5x_1-x_2=9\\-2x_1+3x_2=4\end{cases}$

\item $\begin{cases}x_1+2x_2=8\\3x_1-x_2=3\end{cases}$

\item $\begin{cases}4x_1+3x_2=14\\x_1-2x_2=-1\end{cases}$

\item $\begin{cases}-2x_1+x_2=1\\x_1+5x_2=13\end{cases}$

\item $\begin{cases}7x_1-3x_2=5\\x_1+x_2=3\end{cases}$

\item $\begin{cases}2x_1+x_2=5\\-x_1+4x_2=7\end{cases}$

\item $\begin{cases}x_1-x_2=2\\3x_1+2x_2=11\end{cases}$

\item $\begin{cases}6x_1+x_2=13\\-x_1+3x_2=4\end{cases}$

\item $\begin{cases}-3x_1+2x_2=1\\x_1-x_2=0\end{cases}$

\item $\begin{cases}x_1+x_2=5\\4x_1-3x_2=6\end{cases}$

\item $\begin{cases}2x_1-5x_2=-6\\x_1+x_2=3\end{cases}$

\item $\begin{cases}8x_1+x_2=17\\3x_1-2x_2=0\end{cases}$

\item $\begin{cases}-x_1+2x_2=3\\5x_1-x_2=8\end{cases}$

\item $\begin{cases}3x_1+4x_2=19\\x_1-x_2=1\end{cases}$

\item $\begin{cases}x_1+6x_2=14\\2x_1-x_2=1\end{cases}$

\item $\begin{cases}4x_1-x_2=7\\-x_1+3x_2=2\end{cases}$

\item $\begin{cases}2x_1+3x_2=12\\x_1-2x_2=-1\end{cases}$

\item $\begin{cases}5x_1+2x_2=17\\-x_1+x_2=1\end{cases}$

\item $\begin{cases}x_1-4x_2=-3\\3x_1+x_2=10\end{cases}$

\item $\begin{cases}-2x_1+x_2=-3\\x_1+2x_2=4\end{cases}$

\item $\begin{cases}x_1+x_2=0\\2x_1-3x_2=5\end{cases}$

\item $\begin{cases}9x_1-x_2=8\\x_1+2x_2=5\end{cases}$

\item $\begin{cases}-x_1+5x_2=9\\3x_1+x_2=7\end{cases}$

\item $\begin{cases}2x_1-x_2=3\\-x_1+3x_2=4\end{cases}$

\item $\begin{cases}x_1+7x_2=15\\3x_1-x_2=5\end{cases}$

\item $\begin{cases}4x_1+x_2=10\\x_1-3x_2=-5\end{cases}$

\item $\begin{cases}6x_1-5x_2=2\\x_1+x_2=3\end{cases}$

\end{enumerate}
\end{multicols}

\emph{2. Solve the following $3\times 3$ systems by Gaussian elimination (forward elimination and back substitution).}
\bigskip
\begin{multicols}{2}\raggedright
\begin{enumerate}[label=\arabic*., leftmargin=*]

\item $\begin{cases}x_1+x_2+x_3=6\\x_1-x_2+2x_3=5\\2x_1+x_2-x_3=1\end{cases}$

\item $\begin{cases}2x_1-x_2+3x_3=9\\x_1+2x_2-x_3=0\\3x_1+x_2+2x_3=11\end{cases}$

\item $\begin{cases}x_1+2x_2-x_3=3\\2x_1-x_2+3x_3=1\\x_1+x_2+x_3=6\end{cases}$

\item $\begin{cases}-x_1+x_2+x_3=1\\x_1-x_2+x_3=3\\x_1+x_2-x_3=5\end{cases}$

\item $\begin{cases}3x_1-x_2+2x_3=5\\x_1+2x_2-x_3=0\\2x_1+x_2+3x_3=11\end{cases}$

\item $\begin{cases}x_1+3x_2+2x_3=8\\2x_1-x_2+x_3=3\\-x_1+2x_2-3x_3=-7\end{cases}$

\item $\begin{cases}4x_1-x_2+2x_3=8\\x_1+3x_2-x_3=5\\-2x_1+x_2+3x_3=1\end{cases}$

\item $\begin{cases}x_1-2x_2+3x_3=7\\-x_1+x_2-x_3=-3\\2x_1+3x_2-2x_3=-1\end{cases}$

\item $\begin{cases}2x_1+3x_2-x_3=4\\x_1-x_2+2x_3=7\\3x_1+x_2+x_3=10\end{cases}$

\item $\begin{cases}x_1-x_2+x_3=2\\2x_1+x_2-3x_3=-1\\x_1+2x_2+x_3=7\end{cases}$

\item $\begin{cases}5x_1-2x_2+x_3=3\\x_1+3x_2-2x_3=1\\-x_1+x_2+4x_3=8\end{cases}$

\item $\begin{cases}x_1+x_2-2x_3=-3\\3x_1-2x_2+x_3=4\\2x_1+x_2+3x_3=11\end{cases}$

\item $\begin{cases}-2x_1+x_2+3x_3=4\\x_1+2x_2-x_3=5\\3x_1-x_2+2x_3=6\end{cases}$

\item $\begin{cases}2x_1-3x_2+x_3=0\\x_1+x_2-2x_3=-4\\-x_1+2x_2+x_3=5\end{cases}$

\item $\begin{cases}x_1+4x_2-x_3=7\\-2x_1+x_2+3x_3=2\\x_1-x_2+2x_3=3\end{cases}$

\item $\begin{cases}3x_1-x_2-2x_3=-1\\x_1+2x_2+x_3=8\\-x_1+3x_2-x_3=5\end{cases}$

\item $\begin{cases}x_1-2x_2+x_3=1\\2x_1+x_2-3x_3=0\\-x_1-x_2+2x_3=3\end{cases}$

\item $\begin{cases}4x_1+x_2-x_3=6\\x_1-3x_2+2x_3=-1\\2x_1+2x_2+x_3=9\end{cases}$

\item $\begin{cases}-x_1+2x_2-3x_3=-8\\2x_1-x_2+x_3=5\\x_1+x_2+2x_3=7\end{cases}$

\item $\begin{cases}x_1+x_2+x_3=9\\2x_1+x_2-x_3=2\\x_1-2x_2+3x_3=13\end{cases}$

\item $\begin{cases}3x_1+x_2-2x_3=3\\-x_1+2x_2+x_3=4\\x_1+x_2+3x_3=10\end{cases}$

\item $\begin{cases}2x_1-x_2+x_3=4\\x_1+3x_2-2x_3=-1\\-x_1+x_2+x_3=2\end{cases}$

\item $\begin{cases}x_1-3x_2+2x_3=1\\3x_1+x_2-x_3=8\\2x_1-x_2+3x_3=7\end{cases}$

\item $\begin{cases}-x_1+x_2+4x_3=5\\2x_1+3x_2-x_3=7\\x_1-2x_2+x_3=0\end{cases}$

\item $\begin{cases}x_1+2x_2-3x_3=-4\\-2x_1+x_2+x_3=3\\3x_1-x_2+2x_3=5\end{cases}$

\item $\begin{cases}2x_1+x_2+2x_3=10\\x_1-x_2+x_3=3\\-x_1+2x_2-x_3=0\end{cases}$

\item $\begin{cases}x_1-x_2-x_3=-3\\-x_1+2x_2+3x_3=10\\2x_1+x_2-2x_3=-1\end{cases}$

\item $\begin{cases}3x_1+2x_2-x_3=7\\x_1-x_2+3x_3=6\\-2x_1+x_2+x_3=1\end{cases}$

\item $\begin{cases}x_1+x_2-x_3=4\\-x_1+2x_2+x_3=1\\2x_1-x_2+3x_3=5\end{cases}$

\item $\begin{cases}2x_1+x_2-x_3=3\\x_1-x_2+2x_3=4\\-x_1+2x_2+x_3=5\end{cases}$

\end{enumerate}
 \end{multicols}
\emph{3. Solve the following systems by the Gauss–Jordan method (reduction to RREF).}
\bigskip
\begin{multicols}{2}\raggedright
\begin{enumerate}[label=\arabic*., leftmargin=*]

\item $\begin{cases}x_1+2x_2-x_3=4\\2x_1-x_2+3x_3=1\\x_1+x_2+x_3=6\end{cases}$

\item $\begin{cases}3x_1-x_2=5\\x_1+4x_2=6\end{cases}$

\item $\begin{cases}x_1-x_2+2x_3=3\\2x_1+x_2-x_3=5\\-x_1+3x_2+x_3=4\end{cases}$

\item $\begin{cases}2x_1+x_2=7\\x_1-3x_2=-4\end{cases}$

\item $\begin{cases}-x_1+2x_2+x_3=0\\x_1+x_2-3x_3=4\\2x_1-x_2+x_3=3\end{cases}$

\item $\begin{cases}4x_1-x_2+2x_3=11\\x_1+2x_2-x_3=2\\-2x_1+x_2+3x_3=3\end{cases}$

\item $\begin{cases}x_1+x_2=3\\2x_1-3x_2=-4\end{cases}$

\item $\begin{cases}x_1+x_2+x_3=6\\x_1-x_2+2x_3=5\\2x_1+x_2-x_3=1\end{cases}$

\item $\begin{cases}3x_1+2x_2-x_3=7\\x_1-x_2+3x_3=6\\-2x_1+x_2+x_3=1\end{cases}$

\item $\begin{cases}5x_1-2x_2=8\\x_1+x_2=5\end{cases}$

\item $\begin{cases}x_1-2x_2+x_3=1\\2x_1+x_2-3x_3=0\\-x_1-x_2+2x_3=3\end{cases}$

\item $\begin{cases}x_1+3x_2=10\\4x_1-x_2=1\end{cases}$

\item $\begin{cases}2x_1-x_2+x_3=4\\x_1+3x_2-2x_3=-1\\-x_1+x_2+x_3=2\end{cases}$

\item $\begin{cases}-2x_1+x_2=3\\x_1+x_2=6\end{cases}$

\item $\begin{cases}x_1+x_2-2x_3=-3\\3x_1-2x_2+x_3=4\\2x_1+x_2+3x_3=11\end{cases}$

\item $\begin{cases}x_1-x_2+x_3=2\\2x_1+x_2-3x_3=-1\\x_1+2x_2+x_3=7\end{cases}$

\item $\begin{cases}3x_1+x_2=8\\-x_1+2x_2=1\end{cases}$

\item $\begin{cases}4x_1+x_2-x_3=6\\x_1-3x_2+2x_3=-1\\2x_1+2x_2+x_3=9\end{cases}$

\item $\begin{cases}x_1+x_2+x_3=9\\2x_1+x_2-x_3=2\\x_1-2x_2+3x_3=13\end{cases}$

\item $\begin{cases}2x_1-x_2=3\\x_1+4x_2=9\end{cases}$

\item $\begin{cases}x_1+2x_2-3x_3=-4\\-2x_1+x_2+x_3=3\\3x_1-x_2+2x_3=5\end{cases}$

\item $\begin{cases}-x_1+2x_2-3x_3=-8\\2x_1-x_2+x_3=5\\x_1+x_2+2x_3=7\end{cases}$

\item $\begin{cases}x_1+5x_2=11\\3x_1-x_2=5\end{cases}$

\item $\begin{cases}3x_1-x_2-2x_3=-1\\x_1+2x_2+x_3=8\\-x_1+3x_2-x_3=5\end{cases}$

\item $\begin{cases}x_1-3x_2+2x_3=1\\3x_1+x_2-x_3=8\\2x_1-x_2+3x_3=7\end{cases}$

\item $\begin{cases}6x_1+x_2=13\\-x_1+3x_2=4\end{cases}$

\item $\begin{cases}2x_1+x_2+2x_3=10\\x_1-x_2+x_3=3\\-x_1+2x_2-x_3=0\end{cases}$

\item $\begin{cases}-x_1+x_2+4x_3=5\\2x_1+3x_2-x_3=7\\x_1-2x_2+x_3=0\end{cases}$

\item $\begin{cases}7x_1-x_2=6\\x_1+3x_2=8\end{cases}$

\item $\begin{cases}x_1-x_2-x_3=-3\\-x_1+2x_2+3x_3=10\\2x_1+x_2-2x_3=-1\end{cases}$

\end{enumerate}
 \end{multicols}
\emph{4. Use Gaussian elimination to decide whether each system is consistent; if it is, find its general solution.}
\bigskip
\begin{multicols}{2}\raggedright
\small
\begin{enumerate}[label=\arabic*., leftmargin=*]

\item $\begin{cases}x_1+2x_2-x_3=3\\2x_1+4x_2-2x_3=6\\x_1+x_2+x_3=1\end{cases}$

\item $\begin{cases}x_1-x_2+2x_3=1\\2x_1-2x_2+4x_3=3\end{cases}$

\item $\begin{cases}x_1+x_2+x_3=1\\2x_1+2x_2+2x_3=2\\3x_1+3x_2+3x_3=3\end{cases}$

\item $\begin{cases}x_1+2x_2=5\\2x_1+4x_2=10\\3x_1+6x_2=15\end{cases}$

\item $\begin{cases}2x_1-x_2+3x_3=4\\4x_1-2x_2+6x_3=7\end{cases}$

\item $\begin{cases}x_1-x_2+x_3-x_4=2\\2x_1-2x_2+2x_3-2x_4=4\\x_1+x_2-x_3+x_4=0\end{cases}$

\item $\begin{cases}3x_1-x_2=5\\-6x_1+2x_2=-10\end{cases}$

\item $\begin{cases}x_1+2x_2+3x_3=1\\4x_1+5x_2+6x_3=2\\7x_1+8x_2+9x_3=3\end{cases}$

\item $\begin{cases}x_1-2x_2=3\\-2x_1+4x_2=1\end{cases}$

\item $\begin{cases}x_1+x_2-x_3+2x_4=3\\2x_1+2x_2-2x_3+4x_4=6\\x_1-x_2+x_3=1\end{cases}$

\item $\begin{cases}x_1+3x_2-2x_3=0\\2x_1+6x_2-4x_3=0\end{cases}$

\item $\begin{cases}x_1+x_2=3\\2x_1+2x_2=6\\3x_1+3x_2=10\end{cases}$

\item $\begin{cases}x_1-x_2+2x_3=3\\2x_1-2x_2+4x_3=6\\x_1+x_2-x_3=1\end{cases}$

\item $\begin{cases}2x_1+x_2-x_3=4\\4x_1+2x_2-2x_3=8\\-2x_1-x_2+x_3=-4\end{cases}$

\item $\begin{cases}x_1-x_2+x_3=2\\2x_1-2x_2+2x_3=4\\3x_1-3x_2+3x_3=7\end{cases}$

\item $\begin{cases}x_1+2x_2+x_3+x_4=5\\2x_1+4x_2+2x_3+2x_4=10\end{cases}$

\item $\begin{cases}x_1+x_2=4\\x_1-x_2=2\\2x_1+x_2=7\end{cases}$

\item $\begin{cases}x_1+2x_2-x_3=1\\3x_1+6x_2-3x_3=3\\-x_1-2x_2+x_3=-1\end{cases}$

\item $\begin{cases}5x_1-x_2=10\\-10x_1+2x_2=-15\end{cases}$

\item $\begin{cases}x_1+x_2+x_3+x_4=0\\2x_1-x_2+3x_3=5\\x_1+2x_2-x_3+4x_4=-5\end{cases}$

\item $\begin{cases}x_1+2x_2+3x_3=0\\2x_1+4x_2+6x_3=0\\3x_1+6x_2+9x_3=0\end{cases}$

\item $\begin{cases}x_1-x_2=2\\2x_1-2x_2=4\\3x_1-3x_2=6\end{cases}$

\item $\begin{cases}2x_1+3x_2=7\\4x_1+6x_2=15\end{cases}$

\item $\begin{cases}x_1+x_2-x_3=3\\2x_1+2x_2-2x_3=6\\-x_1-x_2+x_3=-3\end{cases}$

\item $\begin{cases}x_1+2x_2-x_3+3x_4=4\\2x_1+4x_2-2x_3+6x_4=8\\x_1+2x_2-x_3+4x_4=5\end{cases}$

\item $\begin{cases}x_1-3x_2=5\\-2x_1+6x_2=-10\\3x_1-9x_2=15\end{cases}$

\item $\begin{cases}x_1+x_2+x_3=5\\2x_1+2x_2+2x_3=10\\x_1+x_2+x_3=8\end{cases}$

\item $\begin{cases}x_1+2x_2=3\\2x_1+4x_2=6\\-x_1-2x_2=-3\end{cases}$

\item $\begin{cases}3x_1-x_2+2x_3=4\\6x_1-2x_2+4x_3=9\end{cases}$

\item $\begin{cases}x_1+x_2+x_3=1\\2x_1+3x_2-x_3=4\\3x_1+4x_2=5\end{cases}$

\end{enumerate}
\end{multicols}